\section{Details of the dimension-ten certificate}
\label{app:arithmetic}

\subsection{Finite arithmetic}

This subsection gives the exact arithmetic support for the block identities
\eqref{eq:block-symmetry} and the finite-field data in
Section~\ref{sec:certificate}.  The Python 3 program below uses only the
standard library and integer arithmetic.  Its input matrix entries are
those of \eqref{eq:explicit-matrices}, with zero-based row indices in the
code.

The function \texttt{checksym} uses only $i<j$ and column indices $a<b$.
The cases $i=j$ and $a=b$ vanish identically, while exchanging either pair
of indices changes the sign of the same difference.  Thus these strict
loops check every equality in \eqref{eq:block-symmetry}.

The code variable \texttt{R} is the coefficient array $[D\mid H]$, and
\texttt{T} is the reduction of $N$ modulo the current prime.  The program
expands each determinant over the $4!$ permutations, collects coefficients
in the order \eqref{eq:monomial-order}, and performs Gaussian elimination
on $H$ modulo each prime.  The columns of \texttt{T} are the reductions of
the columns in \eqref{eq:N}.

For $f(z)=z^{20}+c_1z^{19}+\cdots+c_{20}$, with $c_0=1$, the program uses
the Newton identities
\[
 k c_k=-\sum_{i=0}^{k-1}c_i\tr(N^{k-i})\qquad(1\leq k\leq20).
\]
Division by $k$ is valid at each of $41,131,139$.  Repeated $p$th powers
modulo $\overline f$ then give the residues of $z^{p^k}$, and the polynomial
Euclidean algorithm gives $d_p(k)$.  All inverses are taken in the stated
prime field.  No floating-point approximation or heuristic search is
involved.

The program asserts the integer block identities and prints the modular
data.  Verification of the modular certificate includes comparing the
printed output with the recorded values below.  The operator, projective, and Galois deductions
are proved in Sections~\ref{sec:geometry} and \ref{sec:certificate}, using
the algebraic facts proved in the next subsection.  The program does not
perform a general elimination or Galois-group computation.

From the paper directory, run
\texttt{python3 verification/exact\_certificate.py} with assertions enabled,
without the \texttt{-O} or \texttt{-OO} optimization options, and with
\texttt{PYTHONOPTIMIZE} unset.

The listing uses an identical build copy of the verification source; its line numbers refer to that source file.

\lstinputlisting[firstnumber=1,caption={Exact verification of the matrix identities and modular certificate.},label={lst:certificate}]{../../verification/exact_certificate.py}

The file \texttt{verification/certificate-output.txt} records the values below; an identical build copy is used for typesetting.
Coefficients of
$\overline f$ are in leading-first order, and the gcd-degree lists are
ordered by $k=1,\ldots,20$.  The first line for each prime gives $p$ and
$\det H\bmod p$.

\lstinputlisting[language={},numbers=none,basicstyle=\ttfamily\footnotesize,caption={Recorded output of the exact verification program.},label={lst:certificate-output}]{../../verification/certificate-output.txt}

\subsection{Algebraic facts}
\label{app:certificate-algebra}

We prove the three general lemmas used in Section~\ref{sec:certificate}.
The first supplies a common zero used to prove the affine quotient nonzero;
the second accommodates rational coefficients with denominators prime to $p$;
the third converts squarefree modular factor degrees into Galois permutations.

\begin{proof}[Proof of Lemma~\ref{lem:six-bilinear}]
Let $X_0,\ldots,X_3,Y_0,\ldots,Y_3$ be indeterminates, put $u_{ij}=X_iY_j$,
and consider the finitely generated complex algebra
\[
 \mathcal R=\C[u_{ij}:0\leq i,j<4]
 \subset\C[X_0,\ldots,X_3,Y_0,\ldots,Y_3].
\]
It is a Noetherian domain.  Its closed complex points are exactly the
matrices of rank at most one.  Indeed, the identities
$u_{ij}u_{k\ell}=u_{i\ell}u_{kj}$ force the values at any point to have
rank at most one, and every such matrix is $xy^{\T}$ and defines an
evaluation of $\mathcal R$.  The maximal ideal
$\mathfrak m=(u_{ij}:0\leq i,j<4)$ corresponds to the zero matrix.

For $1\leq a,b\leq4$, put
$P_{a,b}=(u_{ij}:i\geq a\text{ or }j\geq b)$.  The quotient by $P_{a,b}$
is the domain $\C[X_iY_j:0\leq i<a,\ 0\leq j<b]$.  To justify the
identification, use the monomial basis $X^\alpha Y^\beta$ with
$|\alpha|=|\beta|$ for $\mathcal R$: a monomial containing an excluded
variable factors by an excluded generator, while the remaining monomials
stay distinct.  Hence
\[
 0=P_{4,4}\subsetneq P_{3,4}\subsetneq P_{2,4}\subsetneq P_{1,4}
 \subsetneq P_{1,3}\subsetneq P_{1,2}\subsetneq P_{1,1}
 \subsetneq\mathfrak m
\]
is a strict chain of prime ideals.  In particular,
$\operatorname{ht}\mathfrak m\geq7$.

Let $\ell_1,\ldots,\ell_6\in\mathcal R$ be the linear combinations of
the $u_{ij}$ representing the six bilinear forms, and put
$\mathfrak a=(\ell_1,\ldots,\ell_6)$.  If the only common zero on the
rank-at-most-one matrices were the zero matrix, the Hilbert
Nullstellensatz for the finite-type complex algebra $\mathcal R$ would
give $\sqrt{\mathfrak a}=\mathfrak m$.  The prime $\mathfrak m$ would
then be minimal over an ideal generated by six elements.  Krull's height
theorem bounds the height of such a prime in a Noetherian ring by six,
contradicting the chain above
\cite[Tags 00FV and 0BBZ(1)]{Stacks}.  Thus a nonzero rank-one matrix
$xy^{\T}$ annihilates all six forms.  Both factors are nonzero.
\end{proof}

\begin{proof}[Proof of Lemma~\ref{lem:monic-local-reduction}]
Normalize the valuation by $v_p(p)=1$.  In either pair its values on
nonzero elements are integers, and $R_p/pR_p=\F_p$.  For a nonzero
polynomial $a(z)=\sum_i a_i z^i$ over $\mathbb L$, put
$\nu(a)=\min_{a_i\ne0}v_p(a_i)$.  Scaling nonzero polynomials $a,b$ by
$p^{-\nu(a)}$ and $p^{-\nu(b)}$ makes their coefficients integral with
least valuation zero.  Their reductions are nonzero and have nonzero
product in the domain $\F_p[z]$.  Thus
\[
 \nu(ab)=\nu(a)+\nu(b).
\]
For monic $g,h$ one has $\nu(g),\nu(h)\leq0$, whereas the monic
$P\in R_p[z]$ has $\nu(P)=0$.  Additivity forces
$\nu(g)=\nu(h)=0$, so all coefficients of both factors are integral.
Their leading coefficients remain one after reduction, which preserves
their degrees; a monic constant factor is $1$ and retains degree zero.
Any factorization of $P$ into two positive-degree polynomials over
$\mathbb L$ can be normalized to monic factors, whose reductions would
still have positive degree.  This contradicts irreducibility of the
reduction of $P$.
\end{proof}

\begin{proof}[Proof of Lemma~\ref{lem:local-frobenius}]
Embed the coefficients of $F$ into $\mathbb Z_p$ and put
$s=\operatorname{lcm}(s_1,\ldots,s_h)$.  Choose a monic irreducible
$\overline g\in\F_p[X]$ of degree $s$ and a monic coefficient lift
$g\in\mathbb Z_p[X]$; we use the standard finite-field fact that an
irreducible polynomial of every positive degree exists.  Lemma
\ref{lem:monic-local-reduction} for $(\Q_p,\mathbb Z_p)$ shows that $g$
is irreducible over $\Q_p$.  If $\alpha$ is a root, put
\[
 E=\Q_p(\alpha),\qquad
 \mathcal O=\mathbb Z_p[\alpha]\simeq\mathbb Z_p[X]/(g).
\]
The ring $\mathcal O$ is free of rank $s$ over $\mathbb Z_p$, hence
complete and separated for the $p$-adic topology, and
$\mathcal O/p\mathcal O\simeq\F_{p^s}$.  Every element with nonzero
residue is a unit: lift an inverse modulo $p$ and correct it by a
convergent geometric series.

Every simple residue root of a polynomial in $\mathbb Z_p[z]$ lifts
uniquely to a root in $\mathcal O$.  To see this, suppose $r_n$ is a
root modulo $p^n$ whose derivative has nonzero residue.  The first-order
Taylor congruence gives a unique residue class of $c$ such that
$r_{n+1}=r_n+p^nc$ is a root modulo $p^{n+1}$.  Completeness gives a
limiting root.  If two roots have the same simple residue, the polynomial
difference factorization writes their difference times a unit as zero;
the roots therefore coincide.

Every root of $\overline F$ lies in $\F_{p^s}$ and is simple.  Lifting
them gives $\deg F$ distinct roots of $F$ in $\mathcal O$, hence all
the roots.  Finite fields are perfect, so $\overline\alpha^{\,p}$ is a
simple root of $\overline g$.  Lift it to a root $\beta\in\mathcal O$ of
$g$.  Sending $\alpha$ to $\beta$ defines a $\Q_p$-embedding $E\to E$,
because $g$ is irreducible.  This embedding is an automorphism $\sigma$,
as $E$ is finite-dimensional over $\Q_p$.

Since $\beta\in\mathcal O$, we have $\sigma(\mathcal O)\subseteq\mathcal O$.
The induced residue map fixes $\F_p$ and sends $\overline\alpha$ to
$\overline\alpha^{\,p}$, so it is Frobenius on $\F_{p^s}$.  In a
$\mathbb Z_p$-basis of $\mathcal O$, the matrix of this endomorphism
reduces to an invertible residue matrix.  Its determinant is a $p$-adic
unit, and therefore $\sigma(\mathcal O)=\mathcal O$.

If $\rho_r$ is the lifted root of $F$ with residue $r$, uniqueness of
lifting gives $\sigma(\rho_r)=\rho_{r^p}$.  Thus the permutation of the
lifted roots has exactly the Frobenius cycle lengths on their residues,
namely $s_1,\ldots,s_h$.  The subfield of $E$ generated over $\Q$ by
these roots is a rational splitting field of $F$.  The automorphism
$\sigma$ permutes its generating roots and hence preserves this subfield.
Its restriction is the required rational Galois automorphism.  A
$\Q$-isomorphism transports the statement to any chosen splitting field.
\end{proof}
